\documentclass[12pt,letterpaper]{article}
\usepackage{graphicx} % Required for inserting images
\usepackage{xparse}
\usepackage{xcolor}
\usepackage{bbm}
\usepackage{hyperref}
\usepackage{xurl}
\usepackage{amsmath,mathtools}
\usepackage{amsthm}
\usepackage[margin=1in]{geometry}
\usepackage{tikz}
\usepackage[
    backend=biber, style=oscola, sorting=none,
    doi=false,isbn=false,url=false,eprint=false]{biblatex}
%\usepackage{biblatex}
\usepackage{filecontents}
\addbibresource{references.bib}
\setlength\parindent{0pt}

%\usetikzlibrary{matrix,positioning,calc}
\title{Response to reviewers}
\author{Victor Funes-Leal \and Jared Hutchins}
\date{}

%% BibTeX settings
%\usepackage[square,sort,comma,numbers]{natbib}
%\bibliographystyle{apalike}
%\bibpunct{(}{)}{,}{a}{,}{,}


\begin{document}
% \maketitle

\noindent \section*{Response to Reviewer 1} 
\noindent We thank the reviewer for reading our work and giving thoughtful and constructive comments. We have made several changes that you have suggested, and we detail specifically how we have addressed your comments below.

\subsection*{Comment 1}
{\it
\textbf{Inbreeding as an externality?} The paper argues that the costs of inbreeding are an externality borne by the dairy industry as a whole, but that individual dairy bull breeders do not have an incentive to consider these costs in their individual breeding decisions. The second section of the paper attempts to provide a descriptive conceptual framework for this argument, repeatedly comparing
inbreeding to pesticide resistance management (specifically, Bt crops and refugia).
However, I find the comparison to Bt resistance tenuous and more broadly am not
yet convinced that the costs of inbreeding are an externality at all. Rather than using a descriptive analogy to argue their point, I think AJAE readers would be more convinced by a clear, concise theoretical model: Conceptually, pesticide resistance has been well-established through theoretical modeling  (e.g. \cite{Hueth1974}; \cite{Miranowski1986}; \cite{Cho2024}) to present a type of common pool resource (CPR) problem for farmers, whereby any farmer may use pesticides and draw upon the genetic pool of susceptibility in the pest population (nonexcludability), but in doing so decreases the pool available to all farmers using it in the future due to resistance selection (rivalry). These features are borne out and clearly illustrated in bioeconomic models of pesticide resistance. With dairy cattle inbreeding, the nature and existence of the externality is less clear: Where is nonexcludability (or nonrivalry) present? Based on the information in the paper, high discount rates could provide a simple alternative explanation for increased inbreeding: Breeders may simply favor the short-run trait gains over the long-run risks of inbreeding, even if these discounted future risks are internalized in the market. Of course, dairy industry discount rates may likely be higher than is socially optimal (\cite{Wuepper2023}), but this does not mean there is necessarily an externality present. One might also construe inbreeding as a type of pecuniary externality rather than a ``real'' externality, in which competition and market forces among dairy bull breeders have negative consequences—albeit internalized and market-transmitted—on dairy producers (limited variety of non-inbred bull semen, e.g. in the event of a negative homozygous trait coming to light).\\

Part of the difficulty I have seeing the externality argument lies in the vague
description of the risks of inbreeding and their costs. The main long-run risk the
authors describe in section 2 is an increased ``likelihood of unforeseen genetic traits that could harm the industry''... such as? That is all readers get, and they might be left wondering (as I was) why there’s not an incentive for improvements in GT technologies to better identify and predict these negative traits? A quick literature search (I have no background in dairy science or economics) revealed that such efforts may in fact be underway (\cite{Howard2017}). This vagueness about the costs of inbreeding persists in Section 6 (Cost-Benefit Analysis), which uses an old estimate of the marginal costs of inbreeding (24 USD per 1\% increase in inbreeding), without any explanation of the basis for that figure and how it was calculated (does it depend on a discount rate?). Since it plays a pivotal role in the paper’s economic conclusions, this figure deserves more attention and validation (ideally, this paper would generate an updated estimate of the marginal cost of inbreeding, though I understand that may be beyond the paper’s scope).} \\

\textbf{Reply:} Thank you for this comment. It motivated us to include a conceptual framework in this revision that clarifies what the market frictions are that cause more inbreeding than would be otherwise socially optimal. This is now found in Section 2 of the paper.

Our model is not a fully fleshed out and solved mathematical model, but we believe it clarifies some aspects of the market that can answer your questions. Some of the key aspects of the model are:

\begin{itemize}
\item Dairy breeders compete on providing quality $q$ to the market and the buyers of bulls do not differentiate bulls except on their quality $q$. This leads breeders competing to produce the highest $q$ possible because it allows them to capture a high market share.
\item Dairy breeders can either inbreed ($I$) or outcross ($O$). Inbreeding depletes their genetic diversity $d$ but gives a higher gain in $q$. Conversely, outcrossing improves their own diversity $d$ by using another firm's stock of $d$.
\item Under certain conditions, primarily a high stock of genetic diversity, it is a best response to inbreed (choose $I$) because choosing outcrossing ($O$) gives up market share.
\item The primary market friction is that choosing to outcross improves the firm's own $d$ \textbf{and} the value of choosing $O$ for the other firms. The first benefit is internalized by the firm but the second is not, and so a social planner would choose to outcross more than the individual firm would in equilibrium. The fact that the firm chooses to inbreed more than the social planner means there is a welfare loss, as the industry's level of genetic diversity ($d^{ME}$) will be lower than what the social planner would choose ($d^*$).
\item We characterize the effects of genomic testing as increasing the gains from inbreeding relative to outcrossing. This leads to an equilibrium where genetic diversity is even lower (because of more inbreeding), which may or may not make the welfare loss even larger.
\end{itemize}

Returning to your questions, here is how our model answers them:

\begin{itemize}
\item \textit{With dairy cattle inbreeding, the nature and existence of the externality is less clear: Where is nonexcludability (or nonrivalry) present? } Rather than describe this as an externality, we have framed it as a coordination problem, or a ``race to the bottom'' in terms of quality. Rather than fit the traditional definition of an externality, it fits the profile of a coordination problem where competition incentivizes farmers to inbreed to compete without considering the positive benefits of outcrossing to long-run growth for the industry. The welfare loss is the forgone growth in quality from keeping genetic diversity lower than would be socially optimal.
\item \textit{Based on the information in the paper, high discount rates could provide a simple alternative explanation for increased inbreeding: Breeders may simply favor the short-run trait gains over the long-run risks of inbreeding, even if these discounted future risks are internalized in the
market.} It might be the case that discount rates are not socially optimal, but a welfare loss would still persist even if the discount rates were socially optimal. The welfare loss is caused by firm's not considering the benefits of outcrossing for the industry, and not by not considering the future enough.
\item \textit{Part of the difficulty I have seeing the externality argument lies in the vague description of the risks of inbreeding and their costs. The main long-run risk the authors describe in section 2 is an increased ``likelihood of unforeseen genetic traits that could harm the industry''... such as?} On pg. 4, we discuss the case of one bull that caused 420 million dollars in damages from passing on a recessive traits. Regardless, even if such recessive traits could be 100\% identified, the biggest cost of inbreeding is from reducing genetic diversity which limits the scope of genetic improvement. This is demonstrated in the model that competition keeps the genetic diversity of the industry $d^{ME}$ low which leads to less growth in aggregate quality $\bar{Q}$ over time. We have rewritten the paragraph on pg. 4 to emphasize this: 

\begin{quote}
Excessive inbreeding is an example of a coordination failure in the cattle genetics industry. 
Dairy producers are incentivized to increase the quality of their bulls by either inbreeding or outcrossing. 
Inbreeding captures a larger market share from their competitors by producing more quality, but doing so depletes their own stock of genetic diversity.
By outcrossing and using another company's genetic stock, they can replenish their own genetic diversity but gain less market share because their quality improvement will be lower.
These dynamics incentivize dairy breeders to inbreed more than would be beneficial to the industry because it is costly to outcross when the rest of the industry inbreeds.
Because their own inbreeding makes outcrossing less valuable for firms that may want to use those lines, individual breeders also do not fully consider the costs of their own inbreeding decision.
The result is that inbreeding is done more often in the industry than would be socially optimal.
\end{quote}

\item \textit{This vagueness about the costs of inbreeding persists in Section 6 (Cost-Benefit Analysis), which uses an old estimate of the marginal costs of inbreeding (24 USD per 1\% increase in inbreeding), without any explanation of the basis for that figure and how it was calculated (does it depend on a discount rate?). Since it plays a pivotal role in the paper’s economic conclusions, this figure deserves more attention and validation (ideally, this paper would generate an updated estimate of the marginal cost of inbreeding, though I understand that may be beyond the paper’s scope).} We take this point, and Section 6 has been moved to Appendices A and B and drastically reworked to i) update the costs of inbreeding with a more modern estimate ii) clarify how the cost is calculated and incorporated in the industry iii) contextualize the cost in light of the benefits. Moreover, the true cost (welfare loss) is a counterfactual one and cannot be calculated in this paper (the amount of growth the industry would have had experienced with less inbreeding) without a structural model. We note that this is a limitation of what we are doing in this paper: we can determine how much more inbreeding has happened because of genomic testing, but not what the socially optimal level of inbreeding is. We now note this limitation in the paper on pgs. 28 and 29:
\begin{quote}
More research is needed in this area in agricultural economics, as the coordination problems of dairy cattle breeding have not yet been examined by the field.
Future work could make a significant contribution by expanding our conceptual framework to understand what the socially optimal level of inbreeding is and under what circumstances does a technology like genomic testing actually deepen the welfare losses from inbreeding.
A full cost-benefit analysis of genomic testing is, unfortunately, outside the scope of this paper.
Yet, studying the increase in inbreeding is crucial for ensuring a safe and global resilient dairy sector in the future and ensuring that US dairy cattle breeders continue to sustainably raise the productivity of dairy cattle in the future.

\end{quote}
\end{itemize}

We finally note that we have changed the title of the paper to fully reflect our new emphasis on coordination failure instead of externalities. The new title is ``Competing for Quality: Coordination Failure in Dairy Genetics Markets'', which emphasizes coordination failure as the relevant framework rather than externalities.

\subsection*{Comment 2}
{\it
\textbf{Empirical methods and identification}. Overall, the evidence provided seems pretty clear that GT dramatically increased inbreeding rates. However, I do have some concerns, questions, and suggestions regarding the econometric methods used. The paper adopts a frequently used DID-PSM approach, which is compelling, but their setup makes a subtle argument in the definition of the treatment group, with implications for exactly what is being causally estimated: The authors define
the treatment group of bulls as those coming from lineages of “superstar” founder
bulls from 1998-2004, those whose descendants represented a very high proportion
of all future bulls. They compare these superstar founders to those from matched
controls who were not superstars but had comparable (measured) genetic traits.
They then examine (using DID) how inbreeding rates changed between these two groups after the advent of GT technology. Notably, the treatment group is not defined as those bull lineages in which GT was adopted (the authors do not observe
this information, and argue that GT was so widespread as to make estimation of GT
treatment effects infeasible). So, what do these causal estimates tell us? Assuming
the PSM is successful, this makes selection of trait-matched lineages into superstar status as good as random. Within the subsample of superstar and matched
lineages, the treatment here must technically be defined as: $GT\,(post-2010) \times \text{(randomly selection into)}$ superstar status. But to the extent that it also increased inbreeding among non-superstars, this estimation can’t reveal the aggregate effect of GT on inbreeding. Indeed, Figure 7 appears to show that inbreeding in the matched controls increased more rapidly post-GT (especially starting around 2014). \\

These subtle differences matter for the paper’s conclusions and economic
implications. Yet the BCA in section 6 seems to sweep these subtleties aside- the
superstar-specific GT effect is used to calculate the aggregate, cumulative effect of GT on inbreeding, which is not correct given the definition of the treatment in the econometric model. In particular, it omits the costs of increased inbreeding in nonsuperstars.}\\ 

\textbf{Reply}: Your comment is correct. What this paper identifies is the \textbf{relative} increase in inbreeding in popular lines relative to a group of bulls who are similar in traits but less popular. This is not the same as the aggregate increase in inbreeding relative to a counterfactual where genomic testing did not happen at all. Since every bull has been genomically tested, the only way to identify such an effect would be in a structural model that models this counterfactual. What this paper does is identify an increase in inbreeding consistent with an accelerated competition caused by the genomic testing technology. We explain this now on pg. 28:

\begin{quote}
While there are clearly short-term gains to genomic testing, the actual welfare effects are harder to discern.
As we explore in our conceptual framework, there may still be welfare losses relative to the counterfactual of no genomic testing.
However, our paper does not seek to observe such a counterfactual, instead focusing on the increase in inbreeding for popular lines relative to similar but less popular lines.
Yet, reducing genetic diversity by inbreeding increases genetic gain in the short run at the cost of more steady growth in productivity in the long run.
At its worst, excessive inbreeding leads to a collapse in the population whereby there is no more genetic diversity to continue growing production.
\end{quote}

And on pg. 4:

\begin{quote}
In this paper, we examine whether the introduction of genomic testing increased inbreeding by increasing the returns to inbreeding relative to outcrossing for bulls who were well known in the industry.
Genomic testing significantly increased the returns to inbreeding on these lines because, being the most widely used, the information gain from sequencing these lines was highest.
Genomic selection also made it easier to detect productive genes that they want future offspring to inherit through inbreeding.
Since well-known bulls are most likely to also be high quality, we construct a set of counterfactual bulls that are comparable to the popular genetic families in terms of productivity using propensity score matching.
Using a difference-in-difference approach, we then examine inbreeding rates in popular families relative to our counterfactual families before and after the introduction of genomic testing in 2010.
While the counterfactual of an industry without genomic testing is not observable, we focus on identifying the increase in inbreeding relative to bull lines where genomic testing would increase the returns to inbreeding.
\end{quote}

We have substantially scaled back references to the cost-benefit analysis in the text since our method cannot deliver this calculation. 
In this revision, it gets mentioned only in the introduction and the conclusion, as we moved these cost and benefit calculations to Appendices A and B (formerly Section 6). We now detail the assumptions, formulas used, and values, as well as the respective results in this section. 
The primary results we emphasize are that i) the increases in inbreeding imply costs were borne by the industry and ii) the net benefits were positive.
Our measure of net benefits is the Net Merit index, which is corrected for inbreeding depression.
Thus, positive Net Merit growth implies that the benefits outweigh the costs.
This is now explained on pgs. 27-28:

\begin{quote}
It should be noted, however, that the benefits of genomic testing appear to be much larger than the costs of inbreeding depression in the short term. 
The USDA publishes a profitability index called Net Merit that estimates the profitability contribution of a given dairy bull based on its genetic traits.
Even after correcting for inbreeding depression, we estimate that the index increased by 275 USD per animal relative to what it might have been if genetic traits had not increased after 2010.
The highest per-animal cost of inbreeding in our findings is about a 140 USD per animal increase in inbreeding cost relative to lines that are not popular. 
These costs are difficult to compare since one is relative to a control group and the other is not.\footnote{See Appendices A and B for more details on our calculations.}
Regardless, it appears that the net benefits of genomic testing are positive in the short term.
\end{quote}

\subsection*{Comment 3}
{\it
Trait improvement equation (1): I don’t think equation (1) serves any substantive
purpose in the paper. Economists readers won’t be able to make much of its utility (I wasn’t at least, and I have some knowledge of genetics, though none of dairy bull breeding)—the units in this equation are unclear to me, and it’s not linked to any theoretical model that might shed light on the incentives in the industry in relation to the alleged externality. Nor is it used in any of the econometric estimation (as far as I can tell).} \\

\textbf{Reply}: We have removed all references to the breeder's equation (``equation (1)'') in the text. Instead, we explain the main impacts of genomic testing using words on pages 11:

\begin{quote}
Genomic testing allows dairy breeders to determine the genetic traits of dairy bulls far more quickly than before.
Before it was introduced, it took five years to determine a dairy bull's genetic traits.
Once the bull reached sexual maturity around the age of one year, the bull's semen would be frozen and sent to several farms to produce offspring from the bull.
After about 10 months, the offspring would be born and would take roughly another two years to get to production age. 
After a few years of production, the bull now has production data from its offspring that can be used to calculate its genetic traits, in our model $q$.

Genomic testing gathers the same information in a far shorter amount of time by taking a DNA sample of the bull and correlating its individual genes to traits.
In terms of our model, genomic testing caused an acceleration in the rate of quality gain, that is an increase in $\gamma_I$ and $\gamma_O$.
In the case that they both increase the same amount, there would be no change in $d^{ME}$ since the relative benefits of inbreeding have not changed.
However, genomic testing disproportionately benefits inbreeding because it makes identifying beneficial genes easier.
In particular, it makes identifying which offspring inherited the genes easier and so isolating the beneficial gene is easier than it was before.
Since outcrossing is about expanding the number of genes in the firm's genetic families and not isolating a specific gene, outcrossing benefits far less from genomic testing.
Moreover, this isolation works best when there are fewer genes to test for in the population, and so genomic testing can more efficiently isolate a population which shares less genes (meaning a population with a lower genetic diversity).
\end{quote}


\subsection*{Comment 4}
{\it
Examining Figures 7 and 8, I wonder if we are observing that GT led to earlier
increases in inbreeding in superstars v. matched controls (perhaps due to earlier GT adoption?) with a possible leveling-out starting in 2019 (last year observed in data).

One could speculate that after 2019, the non-superstars’ inbreeding rates eventually caught up with those of superstars, perhaps as GT use further penetrated
down into more marginal lineages. Indeed, one might expect that breeders of the
control bulls—having similar genetic traits but less popular for purely random
reasons—might eventually realize they could use GT to more rapidly improve their
lines, communicate that value to dairy producers, and gain more market share.
}\\

\textbf{Reply}: The Figures show that before the introduction of genomic testing, the average unconditional inbreeding rates of both groups (Figure 7) and average conditional inbreeding rates (Figure 8) did not differ substantially from each other. The plots also show that this phenomenon did not change immediately after 2010 due to adoption lags and biology (one year for a bull to become sexually active, plus another year to father new calves), thereby delaying the increase in traits by a further two years. Our method only works until around 2019, because non-superstar lines die out and the number of animals decreases to near zero, leaving us with an unrepresentative control group; therefore, we cannot assert that the control group's inbreeding rates caught up with those of the treated group.

The comment's economic intuition is correct with respect to trait levels catching up across generations due to competition among animal breeders, who bear the cost of producing a subpar animal via decreased sales. However, given that the control group had fewer and fewer descendants over time, it would not appear to be the case that these lines were ever invested in, but instead were abandoned. We now explain this point in the footnote on pg. 26:

\begin{quote}
    The last period in the sample is 2018 instead of 2020 because we cannot make inferences after 2018 due to a reduction in the number of animals in the control sample; the number of comparable descendants is too small for us to consider them a valid comparison group. This drastic reduction is a consequence of lower fertility rates in the control group; since those sires fathered fewer sons, their sons will beget even fewer grandsons, and after a few years, this group collapses.
\end{quote}

\subsection*{Comment 5}
{\it Line-level clustered standard errors. First, please report the number of lines
(clusters) in the estimation table (table notes are also needed throughout- for *’s, standard errors, reference to corresponding model equations, etc.). Also, one could argue that firm-clustering may be desirable as well: You could do both (as long as lines weren’t fully contained within firms, or vice versa), using multiway clustering!} \\

\textbf{Reply}: Table 2 has been corrected; it includes the number of clusters included in the estimation of the cluster-robust covariance matrices for all four models. We have also included a new table, Table 3, in the appendix that provides for two-way clustered standard errors by sire ID and company ID.

\subsection*{Comment 6}
{\it
BCA in Section 6- discounting? It isn’t clear how the inbreeding USD costs and
genetic trait benefits are aggregated over the time period considered, and if a
discount rate was used in that calculation. Moreover, if the costs of inbreeding
include longer-term risks (see above comment), it isn’t clear how the discount rate
might affect these calculations. Your BCA warrants its own appendix section
detailing these calculations and assumptions for readers.} \\

\textbf{Reply}: You are correct in pointing out that we lacked discounting in our back-of-the-envelope cost calculation. We have also included a more detailed explanation of our estimations. We have thus included a Net Present Value calculation, using a 4.5\% interest rate. Concerning your other comment, we have reduced the scope of our claims, since the information we have at this point is not sufficient to perform a complete BCA. 

\subsection*{Comment 7}
{\it PSM balance assessment \& alternative matching estimators. From my
understanding of the matching literature, raw t-stats and p-values in your balance table are not the best way to assess balance; ‘mean standardized differences’ (MSD) seem to have gained prominence as a way of assessing balance, and I personally much prefer a plot of MSDs, rather than tabular format (see \cite{Greifer2025}). Why did you only consider PSM in your matching strategy rather than alternatives such as nearest-neighbor, or inverse-probability weighting on the propensity score? If they haven’t already done so, I encourage the authors to read through the matching chapter of the Causal Inference Mixtape textbook (\cite{Cunningham2021}). Also, considering that superstar status in founder bulls might be related to firm-level marketing effort, can you include a firm-fixed effect in the propensity-score model? Even better, in my opinion, is a within-cluster(firm) matching estimator, where matches are drawn from the same firm, if available, before drawing matches from other firms (\cite{Arpino2016}).} \\

\textbf{Reply}: Figure 14 in the appendix shows the standardized mean differences for all variables included in the logit model used to estimate the propensity score.

Concerning your second question, our objective for using PSM was to construct a ``synthetic'' control group made up by all descendants of high-quality sires who did not become ``superstars''. Therefore, we do not conduct inferences on those sires, but on their progeny, as far forward as possible. Any other matching method variant would yield the same result, since we are interested only in finding matches on the basis of observable traits, but our inferences are conducted on their descendants.

We have included fixed effects for all companies in the matching model, and the results do not change at all, since the same bulls are selected as alternative sires, and so are their lines. 

\subsection*{Comment 8}
{\it In the main text’s description of the estimation models, some of the notation was hard to follow with incomplete definitions (e.g. $inb$, $yob$). I had to hunt around and deduce indirectly from the text to deduce that $k$ indexes firm and $\alpha_{k}$ was the firm-level fixed effect (there’s also a typo in the appendix, with $\alpha_{j}$ used in one instance where $\alpha_{k}$ should be).} \\

\textbf{Reply}: The notation has been corrected in the main text and appendix to account for this comment. You will find the new text on pg. 22:

\subsection*{Comment 9}
{\it Figure 5: Because of the overlapping shading, I found this figure hard to read. One solution could be to use different line dashing/markers for the different overlaid density plots.} \\

\textbf{Reply}: Figure 5 has been amended and replaced by three plots with different line colors to improve data interpretability.

\subsection*{Comment 10}
{\it In general, I think there are too many citations to the pesticide resistance, area-wide pest control, and invasive species literatures (related to my comment above that I think the comparison here is misplaced). } \\

\textbf{Reply}: In this revision, we have de-emphasized that comparison. It is mentioned in the introduction now on pg. 7 but not in the conceptual framework. From the introduction:

\begin{quote}
However, a related literature in agricultural economics has explored similar coordination problems among farmers controlling pests or disease.
In the short run, farmers are incentivized to forgo the costs of controlling a future pest or disease because of a strong incentive to free-ride off of the efforts of others.
This has been documented in the case of planting Bt-resistant crops (\cite{epanchin2015individual,brown2018voluntary,wechsler2018has,livingston2004managing}) and managing citrus greening \cite{lence2023does}. 
\end{quote}

The idea of citing these papers is to provide agricultural economists a point of reference since almost no literature exists on strategic actions of animal breeders.
We hope that in this revision our conceptual framework will provide enough of a framing that these references only need to be in passing.
The relevant comparison is that the literature on pesticide resistance emphasizes that actors in agriculture might independently in a way that leads to worse outcomes for the whole industry.

%\bibliographystyle{chicago}
%\bibliography{references}
\printbibliography[title={References 1}]

\newrefsection
\newpage
\noindent \section*{Response to Reviewer 2} 
\noindent We thank the reviewer for reviewing our work and giving thoughtful and constructive comments. We have made several changes that you have suggested, and we detail specifically how we have addressed your comments below.

\subsection*{Comment 1}
{\it In the U.S. dairy industry, what is the approximate proportion of commercial dairy breeders who specialize in breeding and sell bull semen as the main business? And what is the proportion of dairy farmers who breed (most of) their cattle and raise them for dairy products? How does the split of the two types of businesses in the authors’ data set compared to that for the entire U.S. industry? Do the two types of business managers
differ significantly in their decisions on breeding (e.g., for the most profitable animal genetics to sell vs. for most productive animals to raise) and use of genomic testing (e.g., to grab market share vs. to avoid inbreeding)?} \\

\textbf{Reply:} In our revision, we have clarified that dairy breeders are completely distinct from dairy farmers. In the dairy industry, it is very rare to see commercial breeders also engaging in dairy farming, as stud farms typically specialize in breeding only. We have clarified this on pg. 8:

\begin{quote}
To describe the dynamics of the dairy industry, we present a simplified theoretical model that describes the coordination problem inherent to dairy cow breeding.
In the modern dairy industry, breeding and selecting new bulls for sale has become almost exclusively the job of dairy breeders.
In contrast, dairy farmers are largely uninvolved with the breeding process and buy their genetics from these companies instead of producing it themselves or selling genetics.
Thus, we model the dairy breeder as a single breeding enterprise and not a farm also engaged in milk production.
\end{quote}

\subsection*{Comment 2}
{\it Prior to the introduction of genomic testing, why were some lines of bulls more popular than others even though they had comparable or same genetic traits (page 15, top line)?

Was the popularity based upon prices, actual productivity (as opposed to the measured genetic traits), branding effects, or something else? After the introduction of genomic testing, did any already popular lines (e.g., super-sires) become less popular and vice versa (because buyers found out some already popular lines now have higher inbreeding rates)? Is it possible for breeders to line breed different lines of cattle since 2008 as some lines ascended to super-sires while other lines became less popular? In other words, can we guarantee there were no crossovers between the treatment group and the control group since 2008?} \\

\textbf{Reply:} One element in the dairy farmer's decision is pedigree. While a bull could be functionally the same, if it is not well known by the industry then it may be seen as a more risky choice. The popular bulls have more descendants and can be assessed using those descendants. While not documented in any literature, we also understand that there are brand effects in the industry given how pedigrees can become popular by winning dairy cow competitions like the World Dairy Expo.

Another key aspect of our new conceptual model postulates that their popularity is precisely what makes them more valuable for inbreeding after genomic testing. This is now explained on pg. 11:

\begin{quote}
Genomic testing allows dairy breeders to determine the genetic traits of dairy bulls far more quickly than before.
Before it was introduced, it took five years to determine a dairy bull's genetic traits.
Once the bull reached sexual maturity around the age of one year, the bull's semen would be frozen and sent to several farms to produce offspring from the bull.
After about 10 months, the offspring would be born and would take roughly another two years to get to production age. 
After a few years of production, the bull now has production data from its offspring that can be used to calculate its genetic traits, in our model $q$.

Genomic testing gathers the same information in a far shorter amount of time by taking a DNA sample of the bull and correlating its individual genes to traits.
In terms of our model, genomic testing caused an acceleration in the rate of quality gain, that is an increase in $\gamma_I$ and $\gamma_O$.
In the case that they both increase the same amount, there would be no change in $d^{ME}$ since the relative benefits of inbreeding have not changed.
However, genomic testing disproportionately benefits inbreeding because it makes identifying beneficial genes easier.
In particular, it makes identifying which offspring inherited the genes easier and so isolating the beneficial gene is easier than it was before.
Since outcrossing is about expanding the number of genes in the firm's genetic families and not isolating a specific gene, outcrossing benefits far less from genomic testing.
Moreover, this isolation works best when there are fewer genes to test for in the population, and so genomic testing can more efficiently isolate a population which shares less genes (meaning a population with a lower genetic diversity).
\end{quote}

As for the point of lines becoming more popular over time, we find that our control lines die out and do not have more descendants. This is why we cannot extend our analysis past 2019. We note this in the footnote on pg. 26:

\begin{quote}
    The last period in the sample is 2018 instead of 2020 because we cannot make inferences after 2018 due to a reduction in the number of animals in the control sample; the number of comparable descendants is too small for us to consider them a valid comparison group. This drastic reduction is a consequence of lower fertility rates in the control group; since those sires fathered fewer sons, their sons will beget even fewer grandsons, and after a few years, this group collapses.
\end{quote}


\subsection*{Comment 3}
{\it Section 2 Conceptual framework examined the two key, sequential decisions facing animal breeders: to line breed or cross breed; to have genomic testing or not. It would be more informative to establish a unified model that explicitly addresses the interdependence between these two decisions with the objective function of market share maximization (say, to breed animals with the highest trait values relative to the market).
In this way, the authors can claim to explore the coordination between individual dairy breeders and the entire industry by analyzing the strategic actions of animal breeders.} \\ 

\textbf{Reply:} Thank you for this comment. Per your suggestion, Section 2 now contains a conceptual model that, while not a formally solved mathematical model, explains the strategic coordination between animal breeders and how they compete on quality. The most important insights from the model are:
\begin{itemize}
\item Dairy breeders compete on providing quality $q$ to the market and the buyers of bulls do not differentiate bulls except on their quality $q$. This leads breeders competing to produce the highest $q$ possible because it allows them to capture a high market share.
\item Dairy breeders can either inbreed ($I$) or outcross ($O$). Inbreeding depletes their genetic diversity $d$ but gives a higher gain in $q$. Conversely, outcrossing improves their own diversity $d$ by using another firm's stock of $d$.
\item Under certain conditions, primarily a high stock of genetic diversity, it is a best response to inbreed (choose $I$) because choosing outcrossing ($O$) gives up market share.
\item The primary market friction is that choosing to outcross improves the firm's own $d$ \textbf{and} the value of choosing $O$ for the other firms. The first benefit is internalized by the firm but the second is not, and so a social planner would choose to outcross more than the individual firm would in equilibrium. The fact that the firm chooses to inbreed more than the social planner means there is a welfare loss, as the industry's level of genetic diversity ($d^{ME}$) will be lower than what the social planner would choose ($d^*$).
\item We characterize the effects of genomic testing as increasing the gains from inbreeding relative to outcrossing. This leads to an equilibrium where genetic diversity is even lower (because of more inbreeding) which may or may not make the welfare loss even larger.
\end{itemize}

We have also renamed the paper to emphasize the importance of the problem being a coordination failure: ``Competing for Quality: Coordination Failure in Dairy Genetics Markets''
\subsection*{Comment 4}
{\it Regarding the cost-benefit analysis in section 6, the authors admitted that this is only a rough, back-of-the-envelope calculation, although there are still potentially substantial benefits of genomic testing to be included. Once animal breeders decide to have their animals genomically tested, they can ascertain the animals’ traits much sooner without
waiting an additional three or four years for the daughter-proven results. Meanwhile, animal breeders need to pay about \$50-100 per animal for testing and possibly cull the undesirable animals at a loss. The adjusted net benefits of genomic testing for the entire industry are likely to be greater than originally estimated, which may render inbreeding a lesser issue in terms of the economic value.} \\

\textbf{Reply:} We agree with your comment that the primary benefits of genomic testing are the ability to ascertain traits quickly (mentioned on pg. 11) and the ability to cull unproductive animals quicker. In our conceptual model, we represent this as increasing the returns to genetic diversity substantially and, most importantly, increasing the returns to inbreeding relative to outcrossing. We discuss this on pg. 12:

\begin{quote}
    The primary prediction of this model is that the introduction of genomic testing increases inbreeding, which potentially exacerbates a welfare loss in the industry. The welfare loss is derived from the fact that breeders excessively inbreed to stay ahead of their competitors, despite the fact that this depletes genetic diversity and leads to less growth in aggregate quality over time. In this paper, our aim is to test one prediction of the model: that genomic testing increases inbreeding. Specifically, we hypothesize that, after genomic testing, the lines that will become the most inbred are lines that were the most popular before genomic testing. Since most bullswould share DNA with these lines, testing their DNA would gain the most information for firms. Taking our model to data, we would then expect inbreeding to increase the most in the popular lines irrespective of quality, as genomic testing will have made inbreeding more beneficial on these lines than before. This would be consistent with the returns to inbreeding increasing relative to outcrossing.
\end{quote}

Since our paper cannot properly do a full cost-benefit analysis of genomic testing, we have rewritten the paper to focus on figuring out how much inbreeding has been increased in popular lines relative to lines where inbreeding would not become more valuable. Our cost-benefit analysis receives a mention in the new Section 6 Discussion but the details of it are now in Appendices A and B. 

In those sections, we have clarified that the returns to genomic testing are seen in the increase in bull quality, measured by the index Net Merit. Since Net Merit already is adjusted for inbreeding, this means that the increase in quality is higher than the losses from inbreeding depression. However, we also emphasize that the costs of genomic testing are from a lower level of diversity that, compared to the social optimum, delivers less growth in the long run. We explain this now on pg. 31:

\begin{quote}
    Published values of Net Merit are already adjusted for inbreeding, the quality increase will be higher than the losses from inbreeding depression. A more accurate estimation of the benefits requires additional information than the one we have, we need information on the expected future inbreeding for all animals in our sample. We do not have the information required to perform such a calculation.
\end{quote}

\subsection*{Comment 5}
{\it I would also suggest adding more justifications when the authors try to project the estimated effects on inbreeding and net merit value using the authors’ sample data into the U.S. dairy industry. Related to my comment \#1 above, the adoption rates of genomic testing may vary between the (commercial) cattle breeders included in this study and cattle producers in the country, so the projected net benefits of genomic testing for the industry will have to be revised accordingly.} \\

\textbf{Reply:} As we explain in the reply to the first comment, commercial cattle breeders and dairy producers are almost completely distinct businesses in the dairy industry. Moreover, the adoption rate of genomic testing is very high in the industry, and the aggregate benefits are evident in the increase in bull quality over time.

\subsection*{Comment 6}
\begin{enumerate}
    \item Page 2, end of paragraph 1: “industr” should be “industry”.

    \textbf{Reply}: the typo was corrected.
    \item Page 7, line 4 from bottom: a period is missing before Melton et al.

    \textbf{Reply}: the typo was corrected.
    \item Page 12, 2nd paragraph, line 3 from the paragraph end: add some numbers to justify ``a substantial portion''.

    \textbf{Reply}: The paragraph now mentions that the NAAB represents 95\% of all livestock genetic companies in the United States.
    \item Page 14, 1st paragraph of section 4, lines 4-5: Why are bulls that are NOT genomically tested more likely to have better quality traits and be more widely used?

    \textbf{Reply}: The assertion is wrong; it was removed from the text.
    \item Page 16, line 4 from page bottom: Figure 5 in the main text is referred to here. Then drop ``in the Appendix''.

    \textbf{Reply}: the typo was corrected.
    \item Page 16, last sentence of the page: perhaps to conduct formal statistical tests on the similarities between the three distributions: matched, treat, and unmatched.

    \textbf{Reply}: A balance table was included, as a response to a previous comment.
    \item Page 24, line 3 below Figure 8: ``percent points'' should be ``percentage points''.
   
    \textbf{Reply}: the typo was corrected.
\end{enumerate}

%\bibliographystyle{chicago}
% \begin{spacing}{3}
%    \bibliography{references}
%\printbibliography[title={References 2}]
%\newrefsection
\newpage

\section*{Response to Reviewer 3}
\noindent We thank the reviewer for reviewing our work and giving thoughtful and constructive comments. We have made several changes that you have suggested, and we detail specifically how we have addressed your comments below.

\subsection*{Comment 1}
\textit{My main concern with your current paper is mostly an inbreeding estimation paper and not an economics paper.  As such, it seems it was targeted to the wrong journal.  That said, you appear to have developed some expertise in the area, so perhaps more economic work in what is an unfilled area in economics will come from future work.}\\

\textbf{Reply:} Thank you for pointing this out. In our last draft, there were parts of the paper that misled readers into thinking this paper was about estimating inbreeding. In fact, the goal of the paper is to analyze the strategic actions of dairy breeders and how genomic testing has changed their decisions. In our new conceptual model we explain the extent to which increased inbreeding rates are a consequence of intense competition among bull breeders to provide the market with the highest quality bulls, since their profits depend on them.

We have also added a new paragraph in page 15 explaining how the inbreeding rate is calculated by the NNAB using extensive progeny data:

\begin{quote}
    The NAAB lists also contain information about lineage, which we use to track descendants of a bull. Every animal born in the USA and Canada is assigned a unique ID which includes information about its breed, country of birth. 
The NAAB data includes these IDs for the bull and both of its parents, allowing us to build a genetic network of relationships from the data. 
Inbreeding is commonly measured using the \textbf{inbreeding coefficient} ($F_{X}$) (\cite{Wright1922}), defined as the probability that two genes in an individual were both descended from a single gene in an ancestor.

\begin{equation}
        F_{X}=\sum_{ca=1}^{K}\left(\dfrac{1}{2}\right)^{n_{1}+n_{2}+1}(1+F_{ca})
\end{equation}

Where $ca$ denotes the common ancestor, $n_{j}$ is the number of generations separating the common ancestor from the sire (dam), and $F_{ca}$ is the inbreeding of the common ancestor. 
This value can be calculated before the birth of an animal, since it depends on the inbreeding coefficients of its parents. The inbreeding coefficient is reported when a newborn animal is registered in its database.
\end{quote}


\subsection*{Comment 2}
{\it Does figure 2 have a source?  Also, what are units on inbreeding rate and how is it calculated? You use this inbreeding term quite a bit and it needs more explanation and definition. Same for TPI on units and how calc in fig 2, though not as important for your paper as inbreeding rate. } \\

\textbf{Reply}: A source was added to Figure 2. Inbreeding is measured with the \textbf{Inbreeding coefficient}, which is equal to the probability that two alleles at the same locus are identical by descent, that is, the probability of inheriting the alleles from their dam and sire. In practice, the CDCB calculates this number using \textbf{progeny} data. Consequently, the inbreeding rate is bounded between zero (absolutely no relation) and 1 (a calf born from parents that are siblings and thus genetically equal to them). \\

The Total Productivity Index is a selection index, a weighted average of PTA values that assigns more weight to specific traits over others to summarize information into a single value per animal, which can be used to rank bulls from highest to lowest. \\

An explanation of both measures has been included in the text to avoid further confusion.
On pg. 14, we explain the inbreeding coefficient and the TPI index, which is explained in the footnote on page 3. \\


\subsection*{Comment 3}
{\it Page 6.  Your back-of-the-envelope cost of inbreeding that you hang your hat on and note in the abstract of your paper is at much at \$6.7 billion. What is this cost? You mention Holstein USA but I do not see this reference in your reference list. When I review the publication at:  \url{https://www.holsteinusa.com/pdf/Interpreting_and_Utilizing_Haplotype_Information_1218.pdf} (which is in your references), I see no mention of the cost you calculate, but rather information about specific DNA sequences that are of concern. Then further, you state that the added milk productivity resulted in a net gain that exceeded this \$6.7 billion cost. So, net, genetic testing to date has been a huge benefit to the dairy industry if I am interpreting your claims correctly. So, why are you anchoring on this “cost” measure of \$6.7 billion.} \\

\textbf{Reply}: Thank you for this comment. Upon reflection, we have realized that this paper cannot do an effective cost-benefit analysis. Instead, the main goal of the paper is to analyze the actions of dairy breeders following the introduction of genomic testing. In this new revision, we have de-emphasized the costs that we calculate here. We incorporate a discount rate and a more detailed explanation of costs under two distinct sets of assumptions. A more detailed appraisal of costs and benefits would require a re-estimation of the Net Merit coefficient to reflect changes in both inbreeding and trait weights, but such a calculation is beyond the scope of this article and merits an independent paper by itself.


\subsection*{Comment 4}
{\it How much does genetic testing cost? Seems we need to know that to also assess net value.} \\

\textbf{Reply}: A single genetic test costs \$35. Source:
\url{https://vgl.ucdavis.edu/test/parentage-genetic-marker-report-cattle}. However, as we explained in the comment above, we no longer frame the cost=benefit analysis as our main contribution.

\subsection*{Comment 5}
{\it Your premise that inbreeding is a problem because it “may have hidden costs” and not that it has been demonstrated to be an economic problem. That is, you rely on conjecture that inbreeding will result in major disease or other genetic concerns that lead I guess to important economic consequences because some traits are being ignored or are not exhibited sufficiently due to inbreeding. While this may be true from a genetics perspective, it seems like conjecture on your part at this point and has no economic assessment in your paper. The paper by Littledon (2024) \url{https://www.rroij.com/open-access/the-role-of-genomics-in-enhancing-dairy-cattle-productivity-and-health.pdf }  indicates genomics testing is enabling selecting for milk yield, feed efficiency, and detection of inherited deficiencies and mutations – essentially a more productive and healthy dairy herd. So, has inbreeding become problematic or do we know when it is a problem? What is the rate of inbreeding at which it is a problem? Are we at a critical level economically yet or are we a long ways away?}  \\

\textbf{Reply:} Thank you for this comment. In this revision, we have emphasized that the costs of inbreeding are from i) inbreeding depression which manifests immediately and ii) lost welfare from keeping the level of genetic diversity lower than would be socially optimal. When there is less genetic diversity, there is less scope for genetic improvement, as having the whole gene pool made up of essentially the same animal would mean there are no more improvements to make.


 This point is now made on pg. 7:

\begin{quote}
If inbreeding continues to increase in the industry, even more damages might occur.
Studies already estimate that the Holstein breed has an unsustainable level of inbreeding, having an ``effective population'' of between 43 and 66; it is commonly assumed that 50 is a critical level under which the breed will lose resilience (\cite{Makanjuola2020}).
Moreover, the lower the genetic diversity, the greater the risk there is of unknown recessive traits causing significant damage to genetic lines (\cite{Cole2011, Lozada2024}). 
At its worst, the population can reach the maximum possible increase in genetic trait values leading to a stalling of genetic improvement.
Future work in economics should seek to understand how policy might be used to address market frictions and incentivize breeders to outcross to ameliorate the impacts of inbreeding.
\end{quote}
 
\subsection*{Comment 6}
{\it Dairy producers are in business for the long run, so why would they act irrationally and select sires for only short-run gains and ignore disease resistance? Do you have any evidence that they actually do this? That is, breeders must provide genetics demanded by producers for long-run decisions. In the Bt corn case you discuss it was the seed companies that required refuge being planted so they could ensure the technology would be reliable in accomplishing what it claimed for the long run. This was not a federal government mandate. So, why would the same not occur in dairy genetics? You imply federal intervention is warranted because of an externality but you have not demonstrated such an externality exists or its magnitude. } \\

\textbf{Reply:} In our new revision, our conceptual model shows that it is individually rational for a dairy breeder to inbreed more than is socially optimal because of the competitive pressures of the industry. Also, outcrossing provides both private and public value to firms but they are only incentivized to care about the private value, which means they will choose outcrossing less than would be optimal to maximize welfare.

\subsection*{Comment 7}
{\it Do dairy producers not also conduct genetic tests increasingly on young heifers for deciding which ones to retain? If so, wouldn’t they also match that genetic info with sires for better complementarity of overall traits?} \\

\textbf{Reply}: Dairy farmers can and do genetically test their female cattle to determine which cow they should mate with each bull on the market. However, it is unclear whether dairy farmers know the inbreeding coefficient of their choices or factor this into their decisions. This points to an information asymmetry issue in the industry. See \cite{Cole2023} for a discussion on how dairy farmers view inbreeding in the industry.

In general, dairy farmers do not keep male cattle within their herds. Most of the semen sold in the market is ``sexed'', meaning that it has a higher probability of giving birth to a female by artificially decreasing the number of Y-chromosomes (see \url{https://www.absglobal.com/uk/what-is-sexed-semen-in-cattle/}). If, by pure chance, a male is born, the farmer is often contractually required to return the animal to the genetics company in exchange for more samples or to sell it for its beef once it is old enough. 

A different breeding strategy involves embryo transfer, in which cows are stimulated to produce a large number of eggs, which are artificially inseminated and transferred to the original cow or a third cow. This strategy also allows for sex-sorting with a very high accuracy. In other words, it is highly unlikely that a dairy farmer keeps bulls in their herds.

\subsection*{Comment 8}
{\it Figure 4. Why has the number of bulls in the market gone up so rapidly from 2010 to 2015?  How does this mesh with your inbreeding measures? Are many of these bulls closely interrelated with each other?  At surface, it would seem more bulls would imply more genetic diversity in the market, but perhaps not?  Might deserve some explanation.} \\

\textbf{Reply:} Thank for this good question. The reason for this increase is that genomic testing allows companies to market bulls before they are sexually mature instead of waiting 5 years for data on their offspring. This allows them to expand their bull offerings. However, this does not necessarily imply that the bulls are not related to each other. In fact, our analysis shows that the average inbreeding coefficient increased after more bulls are introduced, likely because these bulls are even more interrelated than the offerings before genomics. This is itself a consequence of the fact that breeders are more likely to inbreed after genomic testing because it increased the value of inbreeding over outcrossing (see our conceptual model).

\subsection*{Comment 9}
{\it Page 14 you state “Moreover, bulls that are not genomically tested are more likely to have better quality traits and be more widely used.”  Why would bulls that are not tested be more likely to have better quality traits?}

\textbf{Reply}: That sentence was written in error and has been removed from the text.

\subsection*{Comment 10}
{\it Page 16, Figure 5 appears in the text, not the Appendix as stated here.}

\textbf{Reply}: The sentence was corrected to exclude references to the appendix.

%\newpage
%\bibliography{references}
\printbibliography[title={References 3}]
\newrefsection
\end{document}